\item Su equipo ha sido contratado por un importante constructor que está diseñando casas simples para un nuevo fraccionamiento. Él le pide que calcule la cantidad de gas natural que se requerirá para calentar cada casa durante los meses de invierno. La figura TP7.1 muestra la casa en la que está trabajando actualmente. La conductividad térmica promedio de las paredes (incluidas las ventanas) y el techo de la casa representada en la figura es $0.480 \text{ W/m}\cdot^\circ\text{C}$, y su espesor promedio es de $21.0 \text{ cm}$. El calor de la combustión (es decir, la energía suministrada por metro cúbico) de gas natural es $3.89 \times 10^7 \text{ J/m}^3$. 
(a) ¿Cuántos metros cúbicos de gas se deben quemar cada día para mantener una temperatura interior de $25.0 ^\circ\text{C}$ en esta casa si la temperatura exterior es de $0.0 ^\circ\text{C}$? Ignore la radiación y la energía transferida por calor a través del suelo. 
(b) ¿Cómo se verá afectada la respuesta al inciso (a) (\textit{aumentar} o \textit{disminuir} las necesidades de gas) mediante la inclusión de (i) conducción térmica a través del piso; (ii) radiación incidente en el techo, paredes y ventanas durante el día; (iii) operación de electrodomésticos, computadoras, sistemas de entretenimiento; y (iv) fugas de aire a través de grietas alrededor de puertas y ventanas?
